\scanpage{205}
Ainsi, $\pi_1(X)$ n'opère évidemment pas sur $\xi_{x,\mathbf Z}$, ni \unsure{a} sur $\xi_{x,\mathbf R}$, mais l'élément $f_\infty$ de $\pi_1(X)$ (frobénius infini) opère sur $\xi_{x,\mathbf Z}$ et à fortiori sur $\xi_{x,\mathbf R}$ ; on le prolonge \uline{linéairement} en $\tau=f_\infty$ ....

\uline{N.B.} La \uline{conjecture de Hodge} interprète la filtration \uline{induite} sur la coh. \unsure{ordinaire transcendante} $\xi_{x,\mathbf Z}\otimes_{\mathbf Z}\mathbf Q$ \inserted{$=\xi_{x,\mathbf Q}$} par la filtration arithmétique ($=$ motivique, par les niveaux), en termes de la position par. à la graduation [i.e. en termes de la filtration de Hodge -- non motivique -- et du $f_\infty$]. Elle s'exprime en disant que

\deleted{$\unsure{F|E}$ de niveau $\leq\rho\Longleftrightarrow$}

La \uline{partie de niveau} $\leq\rho$ \uline{dans} $\xi_{x,\mathbf Q}(F)$ \uline{n'est autre que}
\[
\xi_{x,\mathbf Q}(F)\cap\xi_{x,\infty}(F)^\rho\cap f_\infty\xi_{x,\infty}(F)^{(\rho)}.
\]

En particulier, le dernier sous-espace de $\xi_{x,\mathbf Q}(F)$ est \uline{motivique} (provient d'un sous-motif $F^{(\rho)}$ de $F$), alors que $\xi_{x,\infty}(F)^{(\rho)}$ et $f_\infty\xi_{x,\infty}(F)^{(\rho)}$ n'est pas motiviques .... \deleted{\unsure{Si $\rho=0$, elle donne} \unsure{l'intersection des réseaux} \unsure{``minimaux''|``universels''}}
\margin{204}
